% ============================================================================
% ARTIGO - SBSI 2026 (Simpósio Brasileiro de Sistemas de Informação)
% Formato: ACM (Association for Computing Machinery)
% Limite: 8-10 páginas (incluindo figuras, tabelas, referências e anexos)
% Revisão: Double-anonymous (anonimizar antes de submeter)
% Idioma: Português ou Inglês
% Evento: 25-28 de maio de 2026, Vitória-ES (presencial)
% Tema: "Sistemas de Informação Inteligentes: Inovações, Aplicações e Ética em IA"
% Site: https://sbsi.sbc.org.br/2026/chamada-pesquisa.php
% ============================================================================

\documentclass[sigconf,nonacm]{acmart}

% ============================================================================
% Configurações ACM
% ============================================================================
\usepackage[utf8]{inputenc}
\usepackage[brazil]{babel}
\usepackage{booktabs}
\usepackage{float}

% Remove ACM copyright/DOI (para submissão)
\setcopyright{none}
\settopmatter{printacmref=false}
\renewcommand\footnotetextcopyrightpermission[1]{}
\pagestyle{plain}

% ============================================================================
% METADADOS
% ============================================================================
\title{Sistema Autônomo de Autenticação Contextual com Blockchain e Inteligência Artificial: Arquitetura de Microserviços para Segurança Adaptativa}

% ============================================================================
% AUTORES
% ATENÇÃO: Para submissão, REMOVER autores (double-anonymous).
% Comente o bloco \author abaixo para versão anônima.
% ============================================================================
\author{João Guedes de Brito}
\email{joao.brito@ifba.edu.br}
\affiliation{%
  \institution{Instituto Federal da Bahia (IFBA)}
  \department{Programa de Pós-Graduação em Engenharia de Sistemas e Produtos}
  \city{Salvador}
  \state{BA}
  \country{Brasil}
}

\author{Cleber Jorge Lira de Santana}
\email{cleber.santana@ifba.edu.br}
\affiliation{%
  \institution{Instituto Federal da Bahia (IFBA)}
  \department{Programa de Pós-Graduação em Engenharia de Sistemas e Produtos}
  \city{Salvador}
  \state{BA}
  \country{Brasil}
}

% ============================================================================
% DOCUMENTO
% ============================================================================
\begin{document}

% ============================================================================
% RESUMO
% ============================================================================
\begin{abstract}
Este artigo apresenta um sistema autônomo de autenticação contextual que integra Inteligência Artificial (IA) e Blockchain em arquitetura de microserviços para segurança adaptativa em aplicações web. O sistema analisa dez atributos contextuais --- geolocalização, dispositivo, horário e padrões comportamentais --- por meio de modelo \textit{ensemble} (Random Forest, Decision Tree, Logistic Regression) para classificação de risco em tempo real. A camada Blockchain garante imutabilidade dos logs de autenticação, atendendo à LGPD. Resultados experimentais demonstram acurácia de 91\% na detecção de anomalias, latência de 320ms e integridade de 100\% nas transações Blockchain. O sistema foi implementado com Spring Boot, Angular, Keycloak, Scikit-learn e Hyperledger/Ethereum, demonstrando viabilidade como alternativa adaptativa aos métodos tradicionais de autenticação.
\end{abstract}

\begin{CCSXML}
<ccs2012>
<concept>
<concept_id>10002978.10002991.10002992</concept_id>
<concept_desc>Security and privacy~Authentication</concept_desc>
<concept_significance>500</concept_significance>
</concept>
<concept>
<concept_id>10010147.10010257</concept_id>
<concept_desc>Computing methodologies~Machine learning</concept_desc>
<concept_significance>300</concept_significance>
</concept>
<concept>
<concept_id>10002978.10003014</concept_id>
<concept_desc>Security and privacy~Distributed systems security</concept_desc>
<concept_significance>300</concept_significance>
</concept>
</ccs2012>
\end{CCSXML}

\ccsdesc[500]{Security and privacy~Authentication}
\ccsdesc[300]{Computing methodologies~Machine learning}
\ccsdesc[300]{Security and privacy~Distributed systems security}

\keywords{Autenticação Contextual, Inteligência Artificial, Blockchain, Machine Learning, Microserviços, LGPD, Segurança da Informação}

\maketitle

% ============================================================================
% 1. INTRODUÇÃO
% ============================================================================
\section{Introdução}

A autenticação de usuários permanece como pilar crítico da segurança da informação em um cenário de crescente sofisticação de ataques cibernéticos. Métodos tradicionais --- senhas fixas e autenticação em dois fatores (2FA) --- apresentam limitações diante de ataques como \textit{credential stuffing} e \textit{session hijacking} \cite{owasp2023}, por não considerarem o contexto comportamental do usuário.

A demanda por conformidade regulatória, como a LGPD \cite{brasil2018} e as diretrizes do NIST \cite{nist2020}, reforça a necessidade de sistemas adaptativos, auditáveis e transparentes. Este trabalho propõe um sistema autônomo de autenticação contextual integrando: (i) análise de risco via Machine Learning; (ii) registro imutável via Blockchain; e (iii) arquitetura de microserviços com práticas DevSecOps.

O sistema analisa dez atributos contextuais por acesso para decisões autônomas sobre concessão, monitoramento, MFA ou bloqueio. A contribuição principal é a integração sistêmica de três camadas de segurança --- identidade (Keycloak/OAuth2), inteligência preditiva (IA) e auditoria imutável (Blockchain) --- em arquitetura modular e reprodutível. Resultados experimentais demonstram acurácia de 91\% e latência compatível com produção.

O tema deste trabalho alinha-se diretamente ao tema central do SBSI 2026 --- ``Sistemas de Informação Inteligentes: Inovações, Aplicações e Ética em IA'' --- ao propor a aplicação de IA em um sistema de informação crítico com considerações éticas de privacidade e conformidade.

% ============================================================================
% 2. TRABALHOS RELACIONADOS
% ============================================================================
\section{Trabalhos Relacionados}

Abordagens de autenticação contextual baseadas em análise comportamental são discutidas por \cite{goodfellow2016}, que demonstram eficácia de modelos supervisionados em dados de alta dimensionalidade. Contudo, a maioria trata análise de risco isoladamente, sem integração com auditoria imutável.

No domínio de Blockchain para segurança, o NISTIR 8202 \cite{nist2018blockchain} descreve a aplicabilidade de blockchains permissionadas para trilhas de auditoria corporativas. O AI Risk Management Framework \cite{nist2023ai} destaca a importância de rastreabilidade e governança em sistemas baseados em IA.

Soluções comerciais (Azure AD Identity Protection, Google BeyondCorp) implementam análise contextual parcial, porém como sistemas proprietários e centralizados, sem transparência nos critérios decisórios. Este trabalho diferencia-se ao propor integração aberta de identidade, IA preditiva e auditoria imutável em arquitetura modular e reprodutível.

% ============================================================================
% 3. ARQUITETURA DO SISTEMA
% ============================================================================
\section{Arquitetura do Sistema}

\subsection{Visão Geral}

A arquitetura organiza-se em cinco camadas: (1) \textbf{Aplicação Modular} --- Spring Boot 3.2.3, Angular 17, Keycloak 24; (2) \textbf{Integração Contínua} --- pipeline DevSecOps automatizado; (3) \textbf{Segurança Automatizada} --- SAST, DAST e escaneamento de dependências; (4) \textbf{Inteligência Artificial} --- modelo ensemble para classificação de risco; (5) \textbf{Registro Imutável} --- Blockchain (Ethereum/Hyperledger Fabric).

O fluxo operacional segue cadeia decisória: tentativa de autenticação $\rightarrow$ captura de contexto $\rightarrow$ estratégias de análise $\rightarrow$ classificação por IA $\rightarrow$ decisão $\rightarrow$ registro Blockchain.

\subsection{Atributos Contextuais}

O sistema captura dez atributos por tentativa de acesso (Tabela~\ref{tab:features}), processados por cinco estratégias de análise (padrão \textit{Strategy}): comportamental, dispositivo, localização, rede e temporal.

\begin{table}[H]
\caption{Atributos contextuais extraídos por acesso}
\label{tab:features}
\small
\begin{tabular}{@{}clp{4.5cm}@{}}
\toprule
\# & Atributo & Descrição \\
\midrule
1  & \texttt{hour\_of\_day}        & Hora do acesso (0--23) \\
2  & \texttt{day\_of\_week}        & Dia da semana (0--6) \\
3  & \texttt{login\_attempts}      & Tentativas na última hora \\
4  & \texttt{is\_new\_device}      & Dispositivo novo (0/1) \\
5  & \texttt{is\_new\_location}    & Localização inédita (0/1) \\
6  & \texttt{is\_vpn}              & Uso de VPN (0/1) \\
7  & \texttt{is\_tor}              & Uso de Tor (0/1) \\
8  & \texttt{session\_duration}    & Duração média de sessão (s) \\
9  & \texttt{pages\_per\_session}  & Páginas por sessão (média) \\
10 & \texttt{time\_since\_login}   & Tempo desde último login (h) \\
\bottomrule
\end{tabular}
\end{table}

\subsection{Modelo de Machine Learning}

O serviço de IA (Python/FastAPI 0.109.0) emprega ensemble com \textit{Soft Voting} combinando Random Forest (100 estimadores, profundidade 10), Decision Tree (profundidade 8) e Logistic Regression (500 iterações). Dados normalizados via \texttt{StandardScaler}; modelo treinado com 5.000 amostras sintéticas.

A classificação segue quatro níveis (Tabela~\ref{tab:risk}):

\begin{table}[H]
\caption{Níveis de risco e decisões automatizadas}
\label{tab:risk}
\small
\begin{tabular}{@{}llll@{}}
\toprule
Nível & Score ($s$) & Decisão & Ação \\
\midrule
Baixo    & $s < 0{,}3$          & PERMITIR  & Acesso liberado \\
Médio    & $[0{,}3; 0{,}6)$     & MONITORAR & Acesso com alerta \\
Alto     & $[0{,}6; 0{,}8)$     & MFA       & Multifator exigido \\
Crítico  & $s \geq 0{,}8$       & BLOQUEAR  & Acesso negado \\
\bottomrule
\end{tabular}
\end{table}

\subsection{Blockchain e Identidade}

A camada Blockchain suporta Ethereum (Web3j 4.10.3) e Hyperledger Fabric 2.5. Cada evento gera transação com: tipo, ID do usuário (anonimizado via \textit{hash}/LGPD), IP, \textit{fingerprint}, score de risco, decisão e timestamp. Integridade verificável por validação criptográfica com confirmação de 12 blocos.

O Keycloak 24 atua como provedor de identidade (OAuth 2.0/OIDC), com proteção contra força bruta, PKCE e tokens de curta duração. Monitoramento via ELK Stack 8.12 e Grafana 10.3.

% ============================================================================
% 4. AVALIAÇÃO EXPERIMENTAL
% ============================================================================
\section{Avaliação Experimental}

\subsection{Ambiente e Metodologia}

Infraestrutura containerizada (Docker) com dez microsserviços: PostgreSQL, Redis, Keycloak, Spring Boot, Angular, FastAPI, Elasticsearch, Logstash, Kibana e Grafana. Três eixos de validação: (i) modelo de ML com validação cruzada ($k=10$); (ii) cadeia completa de autenticação contextual; (iii) integridade Blockchain.

Cenários simulados: acesso habitual; novo dispositivo; VPN/Tor em horário atípico (2h--6h); \textit{credential stuffing}; inatividade prolongada ($>$30 dias).

\subsection{Resultados do Modelo de ML}

\begin{table}[H]
\caption{Métricas de desempenho do classificador}
\label{tab:ml}
\small
\begin{tabular}{@{}ll@{}}
\toprule
Métrica & Valor \\
\midrule
Acurácia & 91,0\% \\
Precisão & 89,4\% \\
Recall   & 93,1\% \\
F1-Score & 91,2\% \\
Falsos Positivos & 8,6\% \\
\bottomrule
\end{tabular}
\end{table}

O recall de 93,1\% prioriza detecção de acessos anômalos, minimizando falsos negativos --- crítico em segurança. A taxa de falsos positivos (8,6\%) é aceitável para validação com dados sintéticos.

\subsection{Latência}

\begin{table}[H]
\caption{Decomposição da latência de autenticação}
\label{tab:latency}
\small
\begin{tabular}{@{}ll@{}}
\toprule
Etapa & Latência \\
\midrule
Keycloak (credenciais)  & 80--100ms \\
Captura de contexto     & 20--30ms \\
Classificação IA        & 120--180ms \\
Registro Blockchain     & 40--60ms \\
\midrule
\textbf{Total}          & \textbf{$\approx$320ms} \\
\bottomrule
\end{tabular}
\end{table}

O módulo de IA é o principal gargalo (120--180ms). O total de 320ms é compatível com padrões de UX (tolerância $\leq$500ms).

\subsection{Integridade Blockchain}

Consistência de 100\% nas transações, com validação criptográfica bem-sucedida em todos os registros. Anonimização por \textit{hash} atende ao princípio de minimização da LGPD.

\subsection{Análise Comparativa}

\begin{table}[H]
\caption{Comparação com métodos tradicionais}
\label{tab:comp}
\small
\begin{tabular}{@{}lccc@{}}
\toprule
Característica & Senha & 2FA & Proposta \\
\midrule
Análise contextual     & Não & Não & Sim \\
Adaptabilidade         & Não & Não & Sim \\
Auditoria imutável     & Não & Não & Sim \\
Detecção de anomalias  & Não & Parcial & Sim \\
Conformidade LGPD      & Parcial & Parcial & Sim \\
Decisão autônoma       & Não & Não & Sim \\
\bottomrule
\end{tabular}
\end{table}

% ============================================================================
% 5. DISCUSSÃO
% ============================================================================
\section{Discussão}

Os resultados confirmam que a integração de IA contextual e Blockchain eleva significativamente a segurança em comparação com métodos estáticos. A abordagem identificou padrões suspeitos em cenários de \textit{credential stuffing} e \textit{session hijacking} invisíveis a sistemas convencionais.

Do ponto de vista ético e de privacidade --- aspecto central do tema do SBSI 2026 --- o sistema implementa anonimização por \textit{hash} antes do registro em Blockchain, minimização de dados pessoais e transparência nos critérios decisórios da IA, alinhando-se às exigências da LGPD e ao AI Risk Management Framework do NIST \cite{nist2023ai}.

As limitações incluem: (i) dados sintéticos para treinamento; (ii) complexidade operacional do Hyperledger Fabric; e (iii) necessidade de políticas refinadas de anonimização. A arquitetura modular permite evolução incremental para endereçar essas limitações.

% ============================================================================
% 6. CONCLUSÃO
% ============================================================================
\section{Conclusão}

Este trabalho apresentou um sistema autônomo de autenticação contextual integrando IA e Blockchain em arquitetura de microserviços. Resultados demonstram acurácia de 91\%, latência de 320ms e integridade total dos registros Blockchain, confirmando viabilidade técnica como alternativa adaptativa aos métodos tradicionais.

A contribuição principal é a integração de três camadas complementares --- identidade, inteligência preditiva e auditoria imutável --- em sistema modular e reprodutível. Trabalhos futuros incluem: dados reais para retreinamento; aprendizado contínuo (\textit{online learning}); testes de carga; extensão para IoT; e explicabilidade (XAI) das decisões do modelo.

% ============================================================================
% REFERÊNCIAS
% ============================================================================
\bibliographystyle{ACM-Reference-Format}

\begin{thebibliography}{10}

\bibitem{brasil2018}
BRASIL. 2018. Lei nº 13.709 -- Lei Geral de Proteção de Dados Pessoais (LGPD). \textit{Diário Oficial da União}, Brasília, DF.

\bibitem{cisa2023}
CISA. 2023. Secure by Design Initiative. Washington, DC.

\bibitem{goodfellow2016}
Ian Goodfellow, Yoshua Bengio, and Aaron Courville. 2016. \textit{Deep Learning}. MIT Press, Cambridge.

\bibitem{nakamoto2008}
Satoshi Nakamoto. 2008. Bitcoin: A Peer-to-Peer Electronic Cash System.

\bibitem{nist2023ai}
NIST. 2023. AI Risk Management Framework (AI RMF 1.0). NIST AI 100-1. Gaithersburg.

\bibitem{nist2018blockchain}
NIST. 2018. Blockchain Technology Overview. NISTIR 8202. Gaithersburg.

\bibitem{nist2022}
NIST. 2022. Secure Software Development Framework (SSDF) v1.1. SP 800-218. Gaithersburg.

\bibitem{nist2020}
NIST. 2020. Security and Privacy Controls for Information Systems and Organizations. SP 800-53 Rev. 5. Gaithersburg.

\bibitem{newman2021}
Sam Newman. 2021. \textit{Building Microservices} (2nd ed.). O'Reilly Media, Sebastopol.

\bibitem{owasp2023}
OWASP Foundation. 2023. OWASP Top 10 -- 2023 Edition.

\bibitem{scikit2024}
Scikit-learn Developers. 2024. Scikit-learn: Machine Learning in Python, v1.4.

\end{thebibliography}

\end{document}
